\documentclass[a4paper,12pt]{scrreprt}

\usepackage{geometry} %% Used for changing geometry of the page
\usepackage[portuguese]{babel} %% Portuguese language
\usepackage[T1]{fontenc} %% Font setup (pdfLaTeX compatible)
\usepackage[utf8]{inputenc}
\usepackage{helvet}
\renewcommand{\familydefault}{\sfdefault} % Arial-like font
\usepackage{xcolor,colortbl} %% Colors and tables
\usepackage{arydshln}
\usepackage[intoc]{nomencl} %% Acronyms
\makenomenclature
\usepackage{graphicx} %% Images
\graphicspath{{images/}}
\usepackage{hyperref} %% Hyperlinks
\hypersetup{
    colorlinks,
    citecolor=black,
    filecolor=black,
    linkcolor=black,
    urlcolor=black
}
\usepackage{amsmath} %% Math
\usepackage[some]{background} %% Background (for cover)
\usepackage{tikz} %% TikZ graphics
\usetikzlibrary{calc}
\usepackage{listings} %% Code listings
\usepackage{float} %% Better float positioning

%% Code listing style
\lstset{
    basicstyle=\ttfamily\footnotesize,
    breaklines=true,
    frame=single,
    language=C,
    showstringspaces=false,
    numbers=left,
    numberstyle=\tiny,
    tabsize=4
}

%% ======================= %%
%% DOCUMENT
\begin{document}
%% Blue cover color
\definecolor{titlepagecolor}{RGB}{120,55,59}

%==========================================================================
% COLORED BAR ON THE LEFT SIDE
%==========================================================================

\backgroundsetup{
    scale=1, 
    angle=0, 
    opacity=1,
    contents={
        \begin{tikzpicture}[remember picture,overlay]
            \path [fill=titlepagecolor] 
                (current page.north west) -- ($(current page.north west) + (5,0)$)
                -- ($(current page.south west) + (5,0)$)-- (current page.south west); 
            \node[color=white] at ($(current page.south west) + (3,4)$) {\bfseries {\fontsize{120}{60} \textsf{LI}}};
            \node[color=titlepagecolor] at ($(current page.south west) + (5.8,4)$) {\bfseries {\fontsize{120}{60} \textsf{ 3}}};
        \end{tikzpicture}
    }
}

%==========================================================================
% TITLE PAGE INFO
%==========================================================================

%% Changes values in this field to show information in the cover and back cover about your team/project


%% TITLE
\title{Relatório Fase 2}

%% AUTHORS
\author{%
\textbf{Maria Fernandes da Silva} \hspace{0.5cm} {\small A106802} \\
\textbf{Matilde Rodrigues Teixeira} \hspace{0.5cm} {\small A106876} \\
\textbf{Francisco Duarte Fernandes Silva} \hspace{0.5cm} {\small A100752}
}

%% Date
\date{\today}

%% Course
\newcommand{\Course}{Licenciatura de Engenharia Informática}

%% Department
\newcommand{\Department}{Escola de Engenharia}

%% UniName
\newcommand{\UniName}{Universidade do Minho}

%% UniPic
\newcommand{\UniPic}{\includegraphics[scale=0.09]{images/uminho.png}}

%% University 
\newcommand{\University}{
    \begin{flushleft}
        \UniPic
    \end{flushleft}
    \textcolor{gray}{\small\textbf{\textsf{\UniName}}}\par
    \textcolor{gray!80!white}{\small{\textsf{\Department}}}\par
    \textcolor{gray!70!white}{\small{\textsf{\Course}}}
}

%% UC
\newcommand{\UC}{
    \begin{flushleft}
        \par\textcolor{titlepagecolor}{  \LARGE\textbf{\textsf{Laboratórios de Informática III}}}
    \end{flushleft}
}

%% School Year
\newcommand{\SchoolYear}{
    \small{\textsf{Ano Letivo de 2025/2026}}}


%% Define new command to show title, author and date
\makeatletter
\let\Title\@title
\let\Author\@author
\let\Date\@date
\makeatother

%==========================================================================
% CLASSIFICATION SECTION 
%==========================================================================

%% School Year
\newcommand{\ReceptionDate}{}
%% Responsible
\newcommand{\Responsible}{}
%% Evaluation
\newcommand{\Evaluation}{}
%% Observations
\newcommand{\Observations}{}





%% MAKETEMPLATE
\newcommand{\makecover}{

%==========================================================================
% BEGIN COVER PAGE 
%==========================================================================

%% Removes page number on footer
\thispagestyle{empty}

%% No indentation 
\setlength{\parindent}{0em}

%% Put Background defined on \backgroundsetup, in this page
\BgThispage

%% Changing geometry to prevent overlay with text
%% At the end of back cover, geometry is default with \restoregeometry
\newgeometry{top=4cm,left=6cm,right=3cm,bottom=2cm}

%% builds university info defined previously
\University
\vspace{1cm}
%% builds curricular unity info defined previously
\UC
%% builds school year info defined previously
\SchoolYear

\vspace*{4cm}
%% bigger space (i think its the default one) between paragraphs 
\setlength{\parskip}{1em}

%% builds title info defined previously
\par\textbf{\textsf{\huge\Title}}
\vspace{1.5cm}
%% builds author(s) info defined previously
\par\Author

\vspace{0.5cm}

%% builds date info defined previously
\par\Date
\restoregeometry
\pagebreak

%==========================================================================
% END COVER PAGE 
%==========================================================================
}
 %% Costume made cover
\pagenumbering{gobble}
\makecover % builds the cover
\newgeometry{top=3cm,left=3cm,right=3cm,bottom=4cm} %% smaller footer and header size

%% BEGIN ABSTRACT PAGE
\renewenvironment{abstract}
 {\par\noindent\textbf{\Large\abstractname}\par\bigskip}
 {}

\pagebreak

%% BEGIN INDEXES PAGES
\renewcommand{\contentsname}{Índice}
\tableofcontents
\pagebreak


%% BEGIN INTRODUCTION
\pagenumbering{arabic}
\chapter{Introdução}

O presente relatório descreve o desenvolvimento da segunda fase do projeto da unidade curricular de Laboratórios de Informática III, implementado em linguagem C. Este projeto centra-se na construção de um sistema destinado à gestão de informação relativa a voos, passageiros, aeronaves, aeroportos e reservas, armazenada inicialmente em ficheiros CSV.

A Fase 2 representa uma evolução significativa face à Fase 1, introduzindo novas funcionalidades e otimizações que visam suportar datasets de maior escala. Nesta fase, foram implementadas três novas queries (Q4, Q5 e Q6) e foi atualizada a Query 1, conforme especificado no enunciado. Adicionalmente, o sistema foi preparado para lidar com um volume de dados substancialmente maior.

Ao longo desta fase, continuaram a ser aplicados os princípios de modularidade e encapsulamento introduzidos na primeira fase, garantindo uma clara separação de responsabilidades entre componentes. Foi dada especial atenção à utilização de estruturas de dados eficientes, nomeadamente hash tables da biblioteca GLib, para assegurar um desempenho adequado nas operações de pesquisa e agregação de dados.

O sistema inclui um programa principal que coordena o fluxo geral de execução, e um programa de testes destinado a validar o comportamento das funcionalidades implementadas. Esta fase teve como objetivo consolidar a base estrutural e funcional do sistema, preparando-o para responder às exigências de desempenho e correção definidas no enunciado.


\chapter{Metodologia utilizada e descrição do sistema}

\section{Arquitetura geral}

A arquitetura do sistema foi concebida de forma modular, mantendo a estrutura base da Fase 1 e expandindo-a para acomodar as novas funcionalidades. O código encontra-se organizado nas seguintes diretorias principais:

\begin{itemize}
    \item \texttt{src/apresentacao/} -- Contém os programas executáveis: \texttt{programa-principal.c} e \texttt{programa-testes.c}
    \item \texttt{src/logica/} -- Módulos de lógica de negócio, incluindo queries, validações e gestão de comandos
    \item \texttt{src/logica/entidades/} -- Definição das entidades do sistema (Airport, Flight, Passenger, Aircraft, Reservation)
    \item \texttt{src/repositorios/} -- Repositórios para gestão das coleções de entidades
    \item \texttt{src/parsers/} -- Módulos de parsing de ficheiros CSV
    \item \texttt{include/} -- Ficheiros de cabeçalho (.h) correspondentes
\end{itemize}

No centro do fluxo encontra-se o programa principal, responsável por iniciar o sistema, carregar os dados através do parser genérico e construir os repositórios das diferentes entidades. Este programa coordena, posteriormente, a execução das queries através do módulo \texttt{gestor-queries.c} e a apresentação dos respetivos resultados.

O módulo \texttt{dados.c} funciona como agregador central, contendo referências para todos os repositórios e estruturas auxiliares necessárias às queries. Esta abordagem permite que as queries acedam aos dados de forma uniforme, sem necessidade de conhecer os detalhes de implementação de cada repositório.

\begin{figure}[H]
\centering
\begin{tikzpicture}[
    box/.style={rectangle, draw, minimum width=2.5cm, minimum height=0.8cm, align=center},
    arrow/.style={->, thick}
]
    % Programa Principal
    \node[box, fill=blue!20] (main) at (0,4) {Programa\\Principal};
    
    % Parser
    \node[box, fill=green!20] (parser) at (-3,2) {Parser};
    
    % Dados
    \node[box, fill=yellow!20] (dados) at (0,2) {Dados};
    
    % Gestor Queries
    \node[box, fill=orange!20] (gestor) at (3,2) {Gestor\\Queries};
    
    % Repositórios
    \node[box, fill=gray!20] (repos) at (-3,0) {Repositórios};
    
    % Entidades
    \node[box, fill=purple!20] (entidades) at (0,0) {Entidades};
    
    % Queries
    \node[box, fill=red!20] (queries) at (3,0) {Q1-Q6};
    
    % Output
    \node[box, fill=cyan!20] (output) at (3,-2) {Output};
    
    % Arrows
    \draw[arrow] (main) -- (parser);
    \draw[arrow] (main) -- (dados);
    \draw[arrow] (main) -- (gestor);
    \draw[arrow] (parser) -- (repos);
    \draw[arrow] (repos) -- (entidades);
    \draw[arrow] (dados) -- (repos);
    \draw[arrow] (gestor) -- (queries);
    \draw[arrow] (queries) -- (dados);
    \draw[arrow] (queries) -- (output);
\end{tikzpicture}
\caption{Arquitetura geral do sistema}
\end{figure}

Esta arquitetura permitiu a construção de um sistema coerente, bem estruturado e facilmente extensível, respeitando os princípios de modularidade e encapsulamento enfatizados nos guiões da unidade curricular.


\section{Estruturas de dados}

Cada entidade é representada por uma \texttt{struct} definida no respetivo ficheiro de implementação, garantindo o encapsulamento. As entidades \texttt{Passenger}, \texttt{Airport}, \texttt{Aircraft}, \texttt{Flight} e \texttt{Reservation} incluem apenas os campos necessários para as validações e para as queries, evitando armazenar informação redundante.

\subsection{Repositórios e Hash Tables}

No que diz respeito ao armazenamento, cada tipo de entidade é guardado no seu repositório através de uma \texttt{GHashTable} da biblioteca GLib. Esta escolha permite pesquisas em tempo constante $O(1)$, uma vez que cada entidade pode ser acedida diretamente através do seu identificador. Por exemplo, o \texttt{RepositorioFlights} mantém:

\begin{itemize}
    \item \texttt{by\_id} -- Hash table principal para acesso por identificador do voo
    \item \texttt{count\_por\_aircraft} -- Hash table para contagem de voos por aeronave (Q2)
    \item \texttt{departures\_por\_data} -- Hash table aninhada para partidas por data e aeroporto (Q3)
    \item \texttt{delay\_por\_companhia} -- Hash table para estatísticas de atrasos por companhia (Q5)
\end{itemize}

\subsection{Estruturas Auxiliares para Queries}

Para operações que exigem ordenação ou percursos sequenciais, são utilizados arrays dinâmicos (\texttt{GArray} ou arrays C convencionais), que facilitam a organização e ordenação dos elementos através da função \texttt{qsort}.

A Query 6 requer uma estrutura auxiliar baseada numa hash table aninhada no módulo \texttt{dados.c}, onde cada chave representa uma nacionalidade e o valor corresponde a outra hash table que mapeia aeroportos para contagens de chegadas.

\subsection{Alocação de Memória}

A alocação de memória é realizada maioritariamente através das funções da GLib (\texttt{g\_new0}, \texttt{g\_slice\_new0}, \texttt{g\_strdup}), que oferecem gestão de memória otimizada. Todas as entidades possuem funções construtoras e destrutoras dedicadas, seguindo o padrão:

\begin{lstlisting}
Flight *flight_new(...);
void flight_free(Flight *f);
\end{lstlisting}


\section{Parsing dos dados}

O parsing foi desenvolvido com base num módulo genérico (\texttt{parser.c}) que pode operar sobre qualquer tipo de entidade, desde que lhe seja fornecida a função de construção apropriada. O parser lê os ficheiros CSV, separa as colunas através da função \texttt{split\_csv\_line}, verifica a validade sintática da linha, converte-a numa entidade se os dados forem válidos e insere a entidade no respetivo repositório.

\subsection{Otimizações de Parsing}

Na Fase 2, foram implementadas várias otimizações no parsing para lidar com datasets de maior dimensão:

\begin{itemize}
    \item \textbf{Buffering aumentado} -- Utilização de \texttt{setvbuf} com buffers de 1 MB tanto para leitura como para escrita de erros
    \item \textbf{Ficheiro de erros único} -- O ficheiro de erros é aberto uma única vez por entidade, em vez de abrir/fechar por linha inválida
    \item \textbf{Função genérica de loading} -- A função \texttt{generic\_loader} centraliza a lógica comum de carregamento, recebendo ponteiros para funções específicas de cada entidade
\end{itemize}

\subsection{Estrutura do Loader Genérico}

O loader genérico recebe como parâmetros:
\begin{itemize}
    \item Diretoria do dataset e nome do ficheiro
    \item Lista de repositórios dependentes (para validações lógicas)
    \item Função de conversão linha $\rightarrow$ objeto
    \item Função de adição ao repositório
    \item Nome do ficheiro de erros
\end{itemize}

Esta abordagem reduziu significativamente a duplicação de código e tornou o sistema mais flexível e modular, alinhando-se com os objetivos pedagógicos da unidade curricular relativos à modularidade.


\chapter{Implementação das funcionalidades}

\section{Validações dos dados}

No projeto, a validação dos dados é realizada em duas fases complementares: a validação sintática e a validação lógica.

\subsection{Validação Sintática}

A validação sintática é efetuada no módulo \texttt{validacoes.c}, onde existem funções específicas para verificar a correção dos campos recebidos do CSV. Nesta fase são analisados aspetos como:

\begin{itemize}
    \item Número de colunas presentes na linha
    \item Formato das datas (\texttt{YYYY-MM-DD} e \texttt{YYYY-MM-DD HH:MM})
    \item Validade de valores numéricos
    \item Inexistência de campos vazios obrigatórios
    \item Formato de emails (username@domain.tld)
    \item Códigos IATA (3 letras maiúsculas)
    \item Identificadores de voos (formato específico)
\end{itemize}

\subsection{Validação Lógica}

Após a verificação sintática, os dados passam por uma segunda fase de validação, de natureza lógica, realizada nos repositórios ou no módulo \texttt{dados.c}. Estas validações dependem da relação entre diferentes entidades e asseguram, por exemplo:

\begin{itemize}
    \item Que um passageiro associado a uma reserva está registado no sistema
    \item Que os voos referenciados numa reserva existem
    \item Que as aeronaves associadas a voos estão registadas
    \item Que as datas de partida e chegada são coerentes
\end{itemize}

Esta distinção entre validações sintáticas e lógicas garante uma separação clara de responsabilidades, reforça a coerência interna do sistema e evita dependências circulares entre módulos.


\section{Queries}

As seis queries implementadas no projeto foram desenvolvidas com base nos requisitos indicados no enunciado, cada uma com objetivos distintos e características próprias ao nível do processamento de dados. Todas se encontram modularizadas nos ficheiros \texttt{query1.c} a \texttt{query6.c}, sendo invocadas pelo gestor de queries após o carregamento integral dos dados.

\subsection{Query 1 -- Informação de Aeroporto}

A Query 1 recebe um código IATA de aeroporto e devolve informação detalhada sobre o mesmo, incluindo nome, cidade, país, tipo, e contagens de chegadas e partidas.

\textbf{Implementação:}
\begin{itemize}
    \item Pesquisa direta na hash table do repositório de aeroportos
    \item As contagens de chegadas e partidas são incrementadas durante o parsing dos voos válidos
    \item Complexidade: $O(1)$ para a pesquisa
\end{itemize}

\subsection{Query 2 -- Top N Aeronaves}

A Query 2 lista as N aeronaves com mais voos realizados, opcionalmente filtradas por fabricante.

\textbf{Implementação:}
\begin{itemize}
    \item Utilização da hash table \texttt{count\_por\_aircraft} pré-computada durante o parsing
    \item Conversão para array e ordenação com \texttt{qsort}
    \item Critério de desempate: ordem alfabética do identificador
    \item Complexidade: $O(n \log n)$ onde $n$ é o número de aeronaves distintas
\end{itemize}

\subsection{Query 3 -- Aeroporto com Mais Partidas}

A Query 3 determina o aeroporto com maior número de partidas num intervalo de datas.

\textbf{Implementação:}
\begin{itemize}
    \item Utilização da estrutura \texttt{departures\_por\_data} (hash table aninhada: data $\rightarrow$ aeroporto $\rightarrow$ contagem)
    \item Iteração sobre as datas no intervalo especificado
    \item Agregação das contagens por aeroporto
    \item Critério de desempate: ordem alfabética do código IATA
\end{itemize}

\subsection{Query 4 -- Passageiro mais tempo no Top 10 (Não Implementada)}

A Query 4 deveria identificar o passageiro que esteve mais tempo no top 10 de passageiros com maior gasto em viagens durante um período.

\textbf{Estado atual:}
\begin{itemize}
    \item A estrutura base foi definida mas a implementação não foi concluída
    \item O ficheiro \texttt{query4.c} contém comentários e estruturas preparatórias
    \item A query está desativada no gestor de queries
\end{itemize}

\subsection{Query 5 -- Top N Companhias com Maior Atraso Médio}

A Query 5 lista as N companhias aéreas com maior atraso médio nos seus voos.

\textbf{Implementação:}
\begin{itemize}
    \item Utilização da hash table \texttt{delay\_por\_companhia} com estatísticas pré-computadas
    \item Estrutura \texttt{AirlineDelayStats} armazena contagem e soma total de atrasos
    \item Cálculo da média: \texttt{total\_delay / count}
    \item Ordenação por média decrescente, desempate alfabético
    \item Precisão de 3 casas decimais no output
\end{itemize}

\subsection{Query 6 -- Aeroporto mais visitado por Nacionalidade}

A Query 6 determina o aeroporto de destino mais visitado por passageiros de uma determinada nacionalidade.

\textbf{Implementação:}
\begin{itemize}
    \item Utilização de hash table aninhada: nacionalidade $\rightarrow$ aeroporto $\rightarrow$ contagem
    \item Normalização da nacionalidade (uppercase, trim)
    \item Critério de desempate: ordem alfabética do código IATA
    \item A estrutura é populada durante a adição de reservas válidas
\end{itemize}


\section{Programa de Testes}

O programa de testes (\texttt{programa-testes.c}) foi desenvolvido para automatizar a validação das queries implementadas. O seu funcionamento consiste em:

\begin{enumerate}
    \item Executar o programa principal com o dataset e ficheiro de input especificados
    \item Comparar os ficheiros de output gerados com os ficheiros esperados
    \item Apresentar estatísticas de correção por query
    \item Reportar discrepâncias encontradas (linha e ficheiro)
    \item Medir tempo de execução e consumo de memória
\end{enumerate}

A distribuição de comandos por query segue o padrão:
\begin{itemize}
    \item Q1: comandos 1-50
    \item Q2: comandos 51-100
    \item Q3: comandos 101-150
    \item Q4: comandos 151-200
    \item Q5: comandos 201-250
    \item Q6: comandos 251-300
\end{itemize}


\chapter{Gestão de memória}

A gestão de memória foi um aspeto central no desenvolvimento do projeto, seguindo as boas práticas enfatizadas nos guiões da unidade curricular, nomeadamente no Guião 5 relativo ao Valgrind.

\section{Estratégias de Alocação}

\subsection{Utilização da GLib}

A biblioteca GLib foi utilizada extensivamente para gestão de memória, oferecendo:

\begin{itemize}
    \item \texttt{g\_new0} e \texttt{g\_slice\_new0} -- Alocação com inicialização a zero
    \item \texttt{g\_strdup} -- Duplicação segura de strings
    \item \texttt{g\_free} e \texttt{g\_slice\_free} -- Libertação correspondente
    \item \texttt{GHashTable} com funções de destruição automática
\end{itemize}

\subsection{Padrão Construtor/Destrutor}

Todas as entidades seguem o padrão de funções construtoras e destrutoras:

\begin{lstlisting}
Passenger *passenger_new(const char *id, ...);
void passenger_free(Passenger *p);
\end{lstlisting}

As hash tables são configuradas com funções de destruição (\texttt{GDestroyNotify}) que libertam automaticamente as chaves e valores quando a tabela é destruída ou quando elementos são removidos.

\section{Resultados do Valgrind}

Na plataforma de testes oficial, foram reportados:

\begin{itemize}
    \item \textbf{Leaks:} 0.195 MB (valor residual aceitável)
    \item \textbf{Memória máxima:} 674 MB (regular) / 4502 MB (large)
\end{itemize}

Os leaks reportados são mínimos e provavelmente relacionados com estruturas internas da GLib que não são libertadas explicitamente antes do término do programa.


\chapter{Análise de Desempenho e Discussão}

\section{Metodologia e Ambiente de Teste}

Os testes foram realizados através da plataforma oficial de testes da unidade curricular, que avalia:
\begin{itemize}
    \item Correção das queries (percentagem de outputs corretos)
    \item Tempo de execução total
    \item Consumo de memória
    \item Existência de memory leaks
\end{itemize}

\section{Resultados Obtidos}

\subsection{Regular Dataset}

\begin{table}[H]
\centering
\begin{tabular}{|l|c|}
\hline
\textbf{Métrica} & \textbf{Valor} \\
\hline
Query 1 & 32\% \\
Query 2 & -- \\
Query 3 & -- \\
Query 4 & 20\% \\
Query 5 & -- \\
Query 6 & 72\% \\
\hline
Tempo Total & 57.6 s \\
Memória & 674 MB \\
Leaks & 0.195 MB \\
\hline
\end{tabular}
\caption{Resultados no Regular Dataset}
\end{table}

\subsection{Large Dataset}

\begin{table}[H]
\centering
\begin{tabular}{|l|c|}
\hline
\textbf{Métrica} & \textbf{Valor} \\
\hline
Query 1 & 32\% \\
Query 2 & -- \\
Query 3 & -- \\
Query 4 & 30\% \\
Query 5 & -- \\
Query 6 & 0\% \\
\hline
Tempo Total & 327.5 s \\
Memória & 4502 MB \\
Leaks & 0.195 MB \\
\hline
\end{tabular}
\caption{Resultados no Large Dataset}
\end{table}

\section{Interpretação dos Resultados}

\subsection{Análise da Query 1}

A Query 1 apresenta uma taxa de correção de 32\% em ambos os datasets. Esta discrepância pode dever-se a:
\begin{itemize}
    \item Problemas na contagem de chegadas/partidas durante o parsing
    \item Diferenças no tratamento de voos cancelados
    \item Possíveis erros na validação de aeroportos
\end{itemize}

\subsection{Análise da Query 4}

A Query 4 não foi completamente implementada, o que explica os resultados baixos (20-30\%). O ficheiro \texttt{query4.c} contém apenas a estrutura base e comentários preparatórios.

\subsection{Análise da Query 6}

A Query 6 apresenta uma discrepância significativa entre datasets:
\begin{itemize}
    \item \textbf{Regular:} 72\% -- Funcionamento parcialmente correto
    \item \textbf{Large:} 0\% -- Falha completa
\end{itemize}

Esta degradação sugere problemas de escalabilidade na implementação, possivelmente relacionados com:
\begin{itemize}
    \item Normalização inconsistente de nacionalidades
    \item Problemas na agregação de dados com volumes maiores
    \item Possíveis race conditions ou overflow em contadores
\end{itemize}

\subsection{Desempenho Temporal}

O tempo de execução escala de forma significativa com o tamanho do dataset:
\begin{itemize}
    \item Regular: 57.6 segundos
    \item Large: 327.5 segundos (aproximadamente 5.7x mais lento)
\end{itemize}

Este fator de crescimento indica que a complexidade temporal não é linear, sugerindo potenciais otimizações nas estruturas de dados ou algoritmos utilizados.

\subsection{Consumo de Memória}

O consumo de memória também escala significativamente:
\begin{itemize}
    \item Regular: 674 MB
    \item Large: 4502 MB (aproximadamente 6.7x mais)
\end{itemize}

Este crescimento é esperado dado o aumento proporcional no volume de dados, embora possam existir oportunidades de otimização.


\chapter{Conclusão}

\section{Objetivos Alcançados}

A segunda fase do projeto permitiu consolidar os conhecimentos adquiridos na primeira fase e aplicá-los num contexto de maior escala. Os principais objetivos alcançados foram:

\begin{itemize}
    \item Manutenção da arquitetura modular estabelecida na Fase 1
    \item Implementação parcial das novas queries (Q4, Q5, Q6)
    \item Utilização de estruturas de dados eficientes (hash tables)
    \item Desenvolvimento de um programa de testes automatizado
    \item Gestão adequada de memória com leaks mínimos
\end{itemize}

\section{Aspetos a Melhorar}

Identificamos os seguintes aspetos que necessitam de melhorias futuras:

\begin{itemize}
    \item \textbf{Query 1:} Rever a lógica de contagem de chegadas/partidas
    \item \textbf{Query 4:} Completar a implementação da análise semanal de gastos
    \item \textbf{Query 6:} Corrigir problemas de escalabilidade no large dataset
    \item \textbf{Desempenho:} Otimizar algoritmos para reduzir tempo de execução
    \item \textbf{Modo Interativo:} Implementar a interface interativa especificada no enunciado
\end{itemize}

\section{Principais Aprendizagens}

O desenvolvimento deste projeto proporcionou aprendizagens valiosas em:

\begin{itemize}
    \item Aplicação prática de modularidade e encapsulamento em C
    \item Utilização de estruturas de dados avançadas (hash tables)
    \item Importância da validação de dados em sistemas reais
    \item Técnicas de debugging com ferramentas como Valgrind e GDB
    \item Gestão de projetos colaborativos com Git
\end{itemize}

\section{Reflexão Crítica}

A segunda fase do projeto revelou-se mais desafiante do que antecipado, particularmente no que diz respeito à implementação das novas queries e à otimização para o large dataset. A Query 4, em particular, requeria uma lógica complexa de análise temporal que não conseguimos completar dentro do prazo.

Os resultados obtidos na plataforma de testes ficaram aquém das expectativas, especialmente nas queries 1, 4 e 6. No entanto, o processo de desenvolvimento permitiu consolidar os conceitos fundamentais de programação em C e preparar-nos para desafios semelhantes em unidades curriculares futuras.

A experiência reforçou a importância de começar a implementação das funcionalidades mais cedo e de realizar testes incrementais ao longo do desenvolvimento, em vez de deixar a validação para o final.


\end{document}
\documentclass[a4paper,12pt]{scrreprt}

\usepackage{geometry} %% Used for changing geometry of the page
\usepackage[portuguese]{babel} %% Portuguese language
\usepackage[T1]{fontenc} %% Font setup (pdfLaTeX compatible)
\usepackage[utf8]{inputenc}
\usepackage{helvet}
\renewcommand{\familydefault}{\sfdefault} % Arial-like font
\usepackage{xcolor,colortbl} %% Colors and tables
\usepackage{arydshln}
\usepackage[intoc]{nomencl} %% Acronyms
\makenomenclature
\usepackage{graphicx} %% Images
\graphicspath{{images/}}
\usepackage{hyperref} %% Hyperlinks
\hypersetup{
    colorlinks,
    citecolor=black,
    filecolor=black,
    linkcolor=black,
    urlcolor=black
}
\usepackage{amsmath} %% Math
\usepackage[some]{background} %% Background (for cover)
\usepackage{tikz} %% TikZ graphics
\usetikzlibrary{calc}
\usepackage{listings} %% Code listings
\usepackage{float} %% Better float positioning

%% Code listing style
\lstset{
    basicstyle=\ttfamily\footnotesize,
    breaklines=true,
    frame=single,
    language=C,
    showstringspaces=false,
    numbers=left,
    numberstyle=\tiny,
    tabsize=4
}

%% ======================= %%
%% DOCUMENT
\begin{document}
%% Blue cover color
\definecolor{titlepagecolor}{RGB}{120,55,59}

%==========================================================================
% COLORED BAR ON THE LEFT SIDE
%==========================================================================

\backgroundsetup{
    scale=1, 
    angle=0, 
    opacity=1,
    contents={
        \begin{tikzpicture}[remember picture,overlay]
            \path [fill=titlepagecolor] 
                (current page.north west) -- ($(current page.north west) + (5,0)$)
                -- ($(current page.south west) + (5,0)$)-- (current page.south west); 
            \node[color=white] at ($(current page.south west) + (3,4)$) {\bfseries {\fontsize{120}{60} \textsf{LI}}};
            \node[color=titlepagecolor] at ($(current page.south west) + (5.8,4)$) {\bfseries {\fontsize{120}{60} \textsf{ 3}}};
        \end{tikzpicture}
    }
}

%==========================================================================
% TITLE PAGE INFO
%==========================================================================

%% Changes values in this field to show information in the cover and back cover about your team/project


%% TITLE
\title{Relatório Fase 2}

%% AUTHORS
\author{%
\textbf{Maria Fernandes da Silva} \hspace{0.5cm} {\small A106802} \\
\textbf{Matilde Rodrigues Teixeira} \hspace{0.5cm} {\small A106876} \\
\textbf{Francisco Duarte Fernandes Silva} \hspace{0.5cm} {\small A100752}
}

%% Date
\date{\today}

%% Course
\newcommand{\Course}{Licenciatura de Engenharia Informática}

%% Department
\newcommand{\Department}{Escola de Engenharia}

%% UniName
\newcommand{\UniName}{Universidade do Minho}

%% UniPic
\newcommand{\UniPic}{\includegraphics[scale=0.09]{images/uminho.png}}

%% University 
\newcommand{\University}{
    \begin{flushleft}
        \UniPic
    \end{flushleft}
    \textcolor{gray}{\small\textbf{\textsf{\UniName}}}\par
    \textcolor{gray!80!white}{\small{\textsf{\Department}}}\par
    \textcolor{gray!70!white}{\small{\textsf{\Course}}}
}

%% UC
\newcommand{\UC}{
    \begin{flushleft}
        \par\textcolor{titlepagecolor}{  \LARGE\textbf{\textsf{Laboratórios de Informática III}}}
    \end{flushleft}
}

%% School Year
\newcommand{\SchoolYear}{
    \small{\textsf{Ano Letivo de 2025/2026}}}


%% Define new command to show title, author and date
\makeatletter
\let\Title\@title
\let\Author\@author
\let\Date\@date
\makeatother

%==========================================================================
% CLASSIFICATION SECTION 
%==========================================================================

%% School Year
\newcommand{\ReceptionDate}{}
%% Responsible
\newcommand{\Responsible}{}
%% Evaluation
\newcommand{\Evaluation}{}
%% Observations
\newcommand{\Observations}{}





%% MAKETEMPLATE
\newcommand{\makecover}{

%==========================================================================
% BEGIN COVER PAGE 
%==========================================================================

%% Removes page number on footer
\thispagestyle{empty}

%% No indentation 
\setlength{\parindent}{0em}

%% Put Background defined on \backgroundsetup, in this page
\BgThispage

%% Changing geometry to prevent overlay with text
%% At the end of back cover, geometry is default with \restoregeometry
\newgeometry{top=4cm,left=6cm,right=3cm,bottom=2cm}

%% builds university info defined previously
\University
\vspace{1cm}
%% builds curricular unity info defined previously
\UC
%% builds school year info defined previously
\SchoolYear

\vspace*{4cm}
%% bigger space (i think its the default one) between paragraphs 
\setlength{\parskip}{1em}

%% builds title info defined previously
\par\textbf{\textsf{\huge\Title}}
\vspace{1.5cm}
%% builds author(s) info defined previously
\par\Author

\vspace{0.5cm}

%% builds date info defined previously
\par\Date
\restoregeometry
\pagebreak

%==========================================================================
% END COVER PAGE 
%==========================================================================
}
 %% Costume made cover
\pagenumbering{gobble}
\makecover % builds the cover
\newgeometry{top=3cm,left=3cm,right=3cm,bottom=4cm} %% smaller footer and header size

%% BEGIN ABSTRACT PAGE
\renewenvironment{abstract}
 {\par\noindent\textbf{\Large\abstractname}\par\bigskip}
 {}

\pagebreak

%% BEGIN INDEXES PAGES
\renewcommand{\contentsname}{Índice}
\tableofcontents
\pagebreak


%% BEGIN INTRODUCTION
\pagenumbering{arabic}
\chapter{Introdução}

O presente relatório descreve o desenvolvimento da segunda fase do projeto da unidade curricular de Laboratórios de Informática III, implementado em linguagem C. Este projeto centra-se na construção de um sistema destinado à gestão de informação relativa a voos, passageiros, aeronaves, aeroportos e reservas, armazenada inicialmente em ficheiros CSV.

A Fase 2 representa uma evolução significativa face à Fase 1, introduzindo novas funcionalidades e otimizações que visam suportar datasets de maior escala. Nesta fase, foram implementadas três novas queries (Q4, Q5 e Q6) e foi atualizada a Query 1, conforme especificado no enunciado. Adicionalmente, o sistema foi preparado para lidar com um volume de dados substancialmente maior.

Ao longo desta fase, continuaram a ser aplicados os princípios de modularidade e encapsulamento introduzidos na primeira fase, garantindo uma clara separação de responsabilidades entre componentes. Foi dada especial atenção à utilização de estruturas de dados eficientes, nomeadamente hash tables da biblioteca GLib, para assegurar um desempenho adequado nas operações de pesquisa e agregação de dados.

O sistema inclui um programa principal que coordena o fluxo geral de execução, e um programa de testes destinado a validar o comportamento das funcionalidades implementadas. Esta fase teve como objetivo consolidar a base estrutural e funcional do sistema, preparando-o para responder às exigências de desempenho e correção definidas no enunciado.


\chapter{Metodologia utilizada e descrição do sistema}

\section{Arquitetura geral}

A arquitetura do sistema foi concebida de forma modular, mantendo a estrutura base da Fase 1 e expandindo-a para acomodar as novas funcionalidades. O código encontra-se organizado nas seguintes diretorias principais:

\begin{itemize}
    \item \texttt{src/apresentacao/} -- Contém os programas executáveis: \texttt{programa-principal.c} e \texttt{programa-testes.c}
    \item \texttt{src/logica/} -- Módulos de lógica de negócio, incluindo queries, validações e gestão de comandos
    \item \texttt{src/logica/entidades/} -- Definição das entidades do sistema (Airport, Flight, Passenger, Aircraft, Reservation)
    \item \texttt{src/repositorios/} -- Repositórios para gestão das coleções de entidades
    \item \texttt{src/parsers/} -- Módulos de parsing de ficheiros CSV
    \item \texttt{include/} -- Ficheiros de cabeçalho (.h) correspondentes
\end{itemize}

No centro do fluxo encontra-se o programa principal, responsável por iniciar o sistema, carregar os dados através do parser genérico e construir os repositórios das diferentes entidades. Este programa coordena, posteriormente, a execução das queries através do módulo \texttt{gestor-queries.c} e a apresentação dos respetivos resultados.

O módulo \texttt{dados.c} funciona como agregador central, contendo referências para todos os repositórios e estruturas auxiliares necessárias às queries. Esta abordagem permite que as queries acedam aos dados de forma uniforme, sem necessidade de conhecer os detalhes de implementação de cada repositório.

\begin{figure}[H]
\centering
\begin{tikzpicture}[
    box/.style={rectangle, draw, minimum width=2.5cm, minimum height=0.8cm, align=center},
    arrow/.style={->, thick}
]
    % Programa Principal
    \node[box, fill=blue!20] (main) at (0,4) {Programa\\Principal};
    
    % Parser
    \node[box, fill=green!20] (parser) at (-3,2) {Parser};
    
    % Dados
    \node[box, fill=yellow!20] (dados) at (0,2) {Dados};
    
    % Gestor Queries
    \node[box, fill=orange!20] (gestor) at (3,2) {Gestor\\Queries};
    
    % Repositórios
    \node[box, fill=gray!20] (repos) at (-3,0) {Repositórios};
    
    % Entidades
    \node[box, fill=purple!20] (entidades) at (0,0) {Entidades};
    
    % Queries
    \node[box, fill=red!20] (queries) at (3,0) {Q1-Q6};
    
    % Output
    \node[box, fill=cyan!20] (output) at (3,-2) {Output};
    
    % Arrows
    \draw[arrow] (main) -- (parser);
    \draw[arrow] (main) -- (dados);
    \draw[arrow] (main) -- (gestor);
    \draw[arrow] (parser) -- (repos);
    \draw[arrow] (repos) -- (entidades);
    \draw[arrow] (dados) -- (repos);
    \draw[arrow] (gestor) -- (queries);
    \draw[arrow] (queries) -- (dados);
    \draw[arrow] (queries) -- (output);
\end{tikzpicture}
\caption{Arquitetura geral do sistema}
\end{figure}

Esta arquitetura permitiu a construção de um sistema coerente, bem estruturado e facilmente extensível, respeitando os princípios de modularidade e encapsulamento enfatizados nos guiões da unidade curricular.


\section{Estruturas de dados}

Cada entidade é representada por uma \texttt{struct} definida no respetivo ficheiro de implementação, garantindo o encapsulamento. As entidades \texttt{Passenger}, \texttt{Airport}, \texttt{Aircraft}, \texttt{Flight} e \texttt{Reservation} incluem apenas os campos necessários para as validações e para as queries, evitando armazenar informação redundante.

\subsection{Repositórios e Hash Tables}

No que diz respeito ao armazenamento, cada tipo de entidade é guardado no seu repositório através de uma \texttt{GHashTable} da biblioteca GLib. Esta escolha permite pesquisas em tempo constante $O(1)$, uma vez que cada entidade pode ser acedida diretamente através do seu identificador. Por exemplo, o \texttt{RepositorioFlights} mantém:

\begin{itemize}
    \item \texttt{by\_id} -- Hash table principal para acesso por identificador do voo
    \item \texttt{count\_por\_aircraft} -- Hash table para contagem de voos por aeronave (Q2)
    \item \texttt{departures\_por\_data} -- Hash table aninhada para partidas por data e aeroporto (Q3)
    \item \texttt{delay\_por\_companhia} -- Hash table para estatísticas de atrasos por companhia (Q5)
\end{itemize}

\subsection{Estruturas Auxiliares para Queries}

Para operações que exigem ordenação ou percursos sequenciais, são utilizados arrays dinâmicos (\texttt{GArray} ou arrays C convencionais), que facilitam a organização e ordenação dos elementos através da função \texttt{qsort}.

A Query 6 requer uma estrutura auxiliar baseada numa hash table aninhada no módulo \texttt{dados.c}, onde cada chave representa uma nacionalidade e o valor corresponde a outra hash table que mapeia aeroportos para contagens de chegadas.

\subsection{Alocação de Memória}

A alocação de memória é realizada maioritariamente através das funções da GLib (\texttt{g\_new0}, \texttt{g\_slice\_new0}, \texttt{g\_strdup}), que oferecem gestão de memória otimizada. Todas as entidades possuem funções construtoras e destrutoras dedicadas, seguindo o padrão:

\begin{lstlisting}
Flight *flight_new(...);
void flight_free(Flight *f);
\end{lstlisting}


\section{Parsing dos dados}

O parsing foi desenvolvido com base num módulo genérico (\texttt{parser.c}) que pode operar sobre qualquer tipo de entidade, desde que lhe seja fornecida a função de construção apropriada. O parser lê os ficheiros CSV, separa as colunas através da função \texttt{split\_csv\_line}, verifica a validade sintática da linha, converte-a numa entidade se os dados forem válidos e insere a entidade no respetivo repositório.

\subsection{Otimizações de Parsing}

Na Fase 2, foram implementadas várias otimizações no parsing para lidar com datasets de maior dimensão:

\begin{itemize}
    \item \textbf{Buffering aumentado} -- Utilização de \texttt{setvbuf} com buffers de 1 MB tanto para leitura como para escrita de erros
    \item \textbf{Ficheiro de erros único} -- O ficheiro de erros é aberto uma única vez por entidade, em vez de abrir/fechar por linha inválida
    \item \textbf{Função genérica de loading} -- A função \texttt{generic\_loader} centraliza a lógica comum de carregamento, recebendo ponteiros para funções específicas de cada entidade
\end{itemize}

\subsection{Estrutura do Loader Genérico}

O loader genérico recebe como parâmetros:
\begin{itemize}
    \item Diretoria do dataset e nome do ficheiro
    \item Lista de repositórios dependentes (para validações lógicas)
    \item Função de conversão linha $\rightarrow$ objeto
    \item Função de adição ao repositório
    \item Nome do ficheiro de erros
\end{itemize}

Esta abordagem reduziu significativamente a duplicação de código e tornou o sistema mais flexível e modular, alinhando-se com os objetivos pedagógicos da unidade curricular relativos à modularidade.


\chapter{Implementação das funcionalidades}

\section{Validações dos dados}

No projeto, a validação dos dados é realizada em duas fases complementares: a validação sintática e a validação lógica.

\subsection{Validação Sintática}

A validação sintática é efetuada no módulo \texttt{validacoes.c}, onde existem funções específicas para verificar a correção dos campos recebidos do CSV. Nesta fase são analisados aspetos como:

\begin{itemize}
    \item Número de colunas presentes na linha
    \item Formato das datas (\texttt{YYYY-MM-DD} e \texttt{YYYY-MM-DD HH:MM})
    \item Validade de valores numéricos
    \item Inexistência de campos vazios obrigatórios
    \item Formato de emails (username@domain.tld)
    \item Códigos IATA (3 letras maiúsculas)
    \item Identificadores de voos (formato específico)
\end{itemize}

\subsection{Validação Lógica}

Após a verificação sintática, os dados passam por uma segunda fase de validação, de natureza lógica, realizada nos repositórios ou no módulo \texttt{dados.c}. Estas validações dependem da relação entre diferentes entidades e asseguram, por exemplo:

\begin{itemize}
    \item Que um passageiro associado a uma reserva está registado no sistema
    \item Que os voos referenciados numa reserva existem
    \item Que as aeronaves associadas a voos estão registadas
    \item Que as datas de partida e chegada são coerentes
\end{itemize}

Esta distinção entre validações sintáticas e lógicas garante uma separação clara de responsabilidades, reforça a coerência interna do sistema e evita dependências circulares entre módulos.


\section{Queries}

As seis queries implementadas no projeto foram desenvolvidas com base nos requisitos indicados no enunciado, cada uma com objetivos distintos e características próprias ao nível do processamento de dados. Todas se encontram modularizadas nos ficheiros \texttt{query1.c} a \texttt{query6.c}, sendo invocadas pelo gestor de queries após o carregamento integral dos dados.

\subsection{Query 1 -- Informação de Aeroporto}

A Query 1 recebe um código IATA de aeroporto e devolve informação detalhada sobre o mesmo, incluindo nome, cidade, país, tipo, e contagens de chegadas e partidas.

\textbf{Implementação:}
\begin{itemize}
    \item Pesquisa direta na hash table do repositório de aeroportos
    \item As contagens de chegadas e partidas são incrementadas durante o parsing dos voos válidos
    \item Complexidade: $O(1)$ para a pesquisa
\end{itemize}

\subsection{Query 2 -- Top N Aeronaves}

A Query 2 lista as N aeronaves com mais voos realizados, opcionalmente filtradas por fabricante.

\textbf{Implementação:}
\begin{itemize}
    \item Utilização da hash table \texttt{count\_por\_aircraft} pré-computada durante o parsing
    \item Conversão para array e ordenação com \texttt{qsort}
    \item Critério de desempate: ordem alfabética do identificador
    \item Complexidade: $O(n \log n)$ onde $n$ é o número de aeronaves distintas
\end{itemize}

\subsection{Query 3 -- Aeroporto com Mais Partidas}

A Query 3 determina o aeroporto com maior número de partidas num intervalo de datas.

\textbf{Implementação:}
\begin{itemize}
    \item Utilização da estrutura \texttt{departures\_por\_data} (hash table aninhada: data $\rightarrow$ aeroporto $\rightarrow$ contagem)
    \item Iteração sobre as datas no intervalo especificado
    \item Agregação das contagens por aeroporto
    \item Critério de desempate: ordem alfabética do código IATA
\end{itemize}

\subsection{Query 4 -- Passageiro mais tempo no Top 10 (Não Implementada)}

A Query 4 deveria identificar o passageiro que esteve mais tempo no top 10 de passageiros com maior gasto em viagens durante um período.

\textbf{Estado atual:}
\begin{itemize}
    \item A estrutura base foi definida mas a implementação não foi concluída
    \item O ficheiro \texttt{query4.c} contém comentários e estruturas preparatórias
    \item A query está desativada no gestor de queries
\end{itemize}

\subsection{Query 5 -- Top N Companhias com Maior Atraso Médio}

A Query 5 lista as N companhias aéreas com maior atraso médio nos seus voos.

\textbf{Implementação:}
\begin{itemize}
    \item Utilização da hash table \texttt{delay\_por\_companhia} com estatísticas pré-computadas
    \item Estrutura \texttt{AirlineDelayStats} armazena contagem e soma total de atrasos
    \item Cálculo da média: \texttt{total\_delay / count}
    \item Ordenação por média decrescente, desempate alfabético
    \item Precisão de 3 casas decimais no output
\end{itemize}

\subsection{Query 6 -- Aeroporto mais visitado por Nacionalidade}

A Query 6 determina o aeroporto de destino mais visitado por passageiros de uma determinada nacionalidade.

\textbf{Implementação:}
\begin{itemize}
    \item Utilização de hash table aninhada: nacionalidade $\rightarrow$ aeroporto $\rightarrow$ contagem
    \item Normalização da nacionalidade (uppercase, trim)
    \item Critério de desempate: ordem alfabética do código IATA
    \item A estrutura é populada durante a adição de reservas válidas
\end{itemize}


\section{Programa de Testes}

O programa de testes (\texttt{programa-testes.c}) foi desenvolvido para automatizar a validação das queries implementadas. O seu funcionamento consiste em:

\begin{enumerate}
    \item Executar o programa principal com o dataset e ficheiro de input especificados
    \item Comparar os ficheiros de output gerados com os ficheiros esperados
    \item Apresentar estatísticas de correção por query
    \item Reportar discrepâncias encontradas (linha e ficheiro)
    \item Medir tempo de execução e consumo de memória
\end{enumerate}

A distribuição de comandos por query segue o padrão:
\begin{itemize}
    \item Q1: comandos 1-50
    \item Q2: comandos 51-100
    \item Q3: comandos 101-150
    \item Q4: comandos 151-200
    \item Q5: comandos 201-250
    \item Q6: comandos 251-300
\end{itemize}


\chapter{Gestão de memória}

A gestão de memória foi um aspeto central no desenvolvimento do projeto, seguindo as boas práticas enfatizadas nos guiões da unidade curricular, nomeadamente no Guião 5 relativo ao Valgrind.

\section{Estratégias de Alocação}

\subsection{Utilização da GLib}

A biblioteca GLib foi utilizada extensivamente para gestão de memória, oferecendo:

\begin{itemize}
    \item \texttt{g\_new0} e \texttt{g\_slice\_new0} -- Alocação com inicialização a zero
    \item \texttt{g\_strdup} -- Duplicação segura de strings
    \item \texttt{g\_free} e \texttt{g\_slice\_free} -- Libertação correspondente
    \item \texttt{GHashTable} com funções de destruição automática
\end{itemize}

\subsection{Padrão Construtor/Destrutor}

Todas as entidades seguem o padrão de funções construtoras e destrutoras:

\begin{lstlisting}
Passenger *passenger_new(const char *id, ...);
void passenger_free(Passenger *p);
\end{lstlisting}

As hash tables são configuradas com funções de destruição (\texttt{GDestroyNotify}) que libertam automaticamente as chaves e valores quando a tabela é destruída ou quando elementos são removidos.

\section{Resultados do Valgrind}

Na plataforma de testes oficial, foram reportados:

\begin{itemize}
    \item \textbf{Leaks:} 0.195 MB (valor residual aceitável)
    \item \textbf{Memória máxima:} 674 MB (regular) / 4502 MB (large)
\end{itemize}

Os leaks reportados são mínimos e provavelmente relacionados com estruturas internas da GLib que não são libertadas explicitamente antes do término do programa.


\chapter{Análise de Desempenho e Discussão}

\section{Metodologia e Ambiente de Teste}

Os testes foram realizados através da plataforma oficial de testes da unidade curricular, que avalia:
\begin{itemize}
    \item Correção das queries (percentagem de outputs corretos)
    \item Tempo de execução total
    \item Consumo de memória
    \item Existência de memory leaks
\end{itemize}

\section{Resultados Obtidos}

\subsection{Regular Dataset}

\begin{table}[H]
\centering
\begin{tabular}{|l|c|}
\hline
\textbf{Métrica} & \textbf{Valor} \\
\hline
Query 1 & 32\% \\
Query 2 & -- \\
Query 3 & -- \\
Query 4 & 20\% \\
Query 5 & -- \\
Query 6 & 72\% \\
\hline
Tempo Total & 57.6 s \\
Memória & 674 MB \\
Leaks & 0.195 MB \\
\hline
\end{tabular}
\caption{Resultados no Regular Dataset}
\end{table}

\subsection{Large Dataset}

\begin{table}[H]
\centering
\begin{tabular}{|l|c|}
\hline
\textbf{Métrica} & \textbf{Valor} \\
\hline
Query 1 & 32\% \\
Query 2 & -- \\
Query 3 & -- \\
Query 4 & 30\% \\
Query 5 & -- \\
Query 6 & 0\% \\
\hline
Tempo Total & 327.5 s \\
Memória & 4502 MB \\
Leaks & 0.195 MB \\
\hline
\end{tabular}
\caption{Resultados no Large Dataset}
\end{table}

\section{Interpretação dos Resultados}

\subsection{Análise da Query 1}

A Query 1 apresenta uma taxa de correção de 32\% em ambos os datasets. Esta discrepância pode dever-se a:
\begin{itemize}
    \item Problemas na contagem de chegadas/partidas durante o parsing
    \item Diferenças no tratamento de voos cancelados
    \item Possíveis erros na validação de aeroportos
\end{itemize}

\subsection{Análise da Query 4}

A Query 4 não foi completamente implementada, o que explica os resultados baixos (20-30\%). O ficheiro \texttt{query4.c} contém apenas a estrutura base e comentários preparatórios.

\subsection{Análise da Query 6}

A Query 6 apresenta uma discrepância significativa entre datasets:
\begin{itemize}
    \item \textbf{Regular:} 72\% -- Funcionamento parcialmente correto
    \item \textbf{Large:} 0\% -- Falha completa
\end{itemize}

Esta degradação sugere problemas de escalabilidade na implementação, possivelmente relacionados com:
\begin{itemize}
    \item Normalização inconsistente de nacionalidades
    \item Problemas na agregação de dados com volumes maiores
    \item Possíveis race conditions ou overflow em contadores
\end{itemize}

\subsection{Desempenho Temporal}

O tempo de execução escala de forma significativa com o tamanho do dataset:
\begin{itemize}
    \item Regular: 57.6 segundos
    \item Large: 327.5 segundos (aproximadamente 5.7x mais lento)
\end{itemize}

Este fator de crescimento indica que a complexidade temporal não é linear, sugerindo potenciais otimizações nas estruturas de dados ou algoritmos utilizados.

\subsection{Consumo de Memória}

O consumo de memória também escala significativamente:
\begin{itemize}
    \item Regular: 674 MB
    \item Large: 4502 MB (aproximadamente 6.7x mais)
\end{itemize}

Este crescimento é esperado dado o aumento proporcional no volume de dados, embora possam existir oportunidades de otimização.


\chapter{Conclusão}

\section{Objetivos Alcançados}

A segunda fase do projeto permitiu consolidar os conhecimentos adquiridos na primeira fase e aplicá-los num contexto de maior escala. Os principais objetivos alcançados foram:

\begin{itemize}
    \item Manutenção da arquitetura modular estabelecida na Fase 1
    \item Implementação parcial das novas queries (Q4, Q5, Q6)
    \item Utilização de estruturas de dados eficientes (hash tables)
    \item Desenvolvimento de um programa de testes automatizado
    \item Gestão adequada de memória com leaks mínimos
\end{itemize}

\section{Aspetos a Melhorar}

Identificamos os seguintes aspetos que necessitam de melhorias futuras:

\begin{itemize}
    \item \textbf{Query 1:} Rever a lógica de contagem de chegadas/partidas
    \item \textbf{Query 4:} Completar a implementação da análise semanal de gastos
    \item \textbf{Query 6:} Corrigir problemas de escalabilidade no large dataset
    \item \textbf{Desempenho:} Otimizar algoritmos para reduzir tempo de execução
    \item \textbf{Modo Interativo:} Implementar a interface interativa especificada no enunciado
\end{itemize}

\section{Principais Aprendizagens}

O desenvolvimento deste projeto proporcionou aprendizagens valiosas em:

\begin{itemize}
    \item Aplicação prática de modularidade e encapsulamento em C
    \item Utilização de estruturas de dados avançadas (hash tables)
    \item Importância da validação de dados em sistemas reais
    \item Técnicas de debugging com ferramentas como Valgrind e GDB
    \item Gestão de projetos colaborativos com Git
\end{itemize}

\section{Reflexão Crítica}

A segunda fase do projeto revelou-se mais desafiante do que antecipado, particularmente no que diz respeito à implementação das novas queries e à otimização para o large dataset. A Query 4, em particular, requeria uma lógica complexa de análise temporal que não conseguimos completar dentro do prazo.

Os resultados obtidos na plataforma de testes ficaram aquém das expectativas, especialmente nas queries 1, 4 e 6. No entanto, o processo de desenvolvimento permitiu consolidar os conceitos fundamentais de programação em C e preparar-nos para desafios semelhantes em unidades curriculares futuras.

A experiência reforçou a importância de começar a implementação das funcionalidades mais cedo e de realizar testes incrementais ao longo do desenvolvimento, em vez de deixar a validação para o final.


\end{document}
